%------------------------------------------------------------------------------%
\begin{question}
  Dada a equação da hipérbole
  \[
    \frac{(x-1)^2}{16} - \frac{(y+2)^2}{9} = 1
  \]
  \vspace{-1.5\baselineskip}
  \begin{enumerate}
    \item encontre o centro e os vértices
    \item escreva a equação matricial
    \item esboce a curva
  \end{enumerate}
\end{question}

%------------------------------------------------------------------------------%
\begin{resolution}{Centro}
  Comparando
  \[
    \frac{(x-1)^2}{16} - \frac{(y+2)^2}{9} = 1
  \]
  com a equação canônica transladada
  \[
    \frac{(x-h)^2}{a^2} - \frac{(y-k)^2}{b^2} = 1
  \]
  obtemos
  \[
    a^2 = 16
    \qquad
    a = 4
  \]
  \[
    b^2 = 9
    \qquad
    b = 3
  \]
  o centro é
  \(
    C
    = (h, k)
    = (1, -2)
  \)
\end{resolution}

%------------------------------------------------------------------------------%
\begin{resolution}{Vértices}
\begin{columns}[t]
\column{0.6\textwidth}
  Novas coordenadas
  \[
    V
    = \begin{bmatrix}
      x - h  \\
      y - k
    \end{bmatrix}
    = \begin{bmatrix}[c]
      x - 2  \\
      y - (-1)
    \end{bmatrix}
    = \begin{bmatrix}
      x - 2  \\
      y + 1
    \end{bmatrix}
  \]
\column{0.4\textwidth}
  Vértices
  \[
    V
    = \begin{bmatrix}
      x - 2  \\
      y + 1
    \end{bmatrix}
    = \begin{bmatrix}
      \pm 4  \\
      0
    \end{bmatrix}
  \]
\end{columns}

\medskip
Resolvendo
\vspace{-1.5\baselineskip}
\begin{columns}[t]
\column{0.2\textwidth}
  \begin{gather*}
    y + 1 = 0 \\
    y     = -1
  \end{gather*}
\column{0.2\textwidth}
  \begin{gather*}
    x - 2 = 4     \\
    x     = 4 + 2 \\
    x     = 6
  \end{gather*}
\column{0.2\textwidth}
  \begin{gather*}
    x - 2 = -4 \\
    x     = -4 + 2 \\
    x     = -2
  \end{gather*}
\end{columns}

  \medskip
  Vértices:
  \(
    \quad
    \mathcal{V}_1 = (-2, -1)
    \quad
    \mathcal{V}_1 = (6, -1)
  \)
\end{resolution}

%------------------------------------------------------------------------------%
\begin{resolution}{Equação Matricial}
  Equação Matricial
  \[
    V^t A V = 1
  \]
  com
  \[
    A = \begin{bmatrix}
      \dfrac{1}{16} & 0            \\[2mm]
      0             & -\dfrac{1}{9}
    \end{bmatrix}
    \qquad\qquad
    V = \begin{bmatrix}
      x - 2  \\
      y + 1
    \end{bmatrix}
  \]
\end{resolution}

%------------------------------------------------------------------------------%
\begin{resolution}{Esboço}
\end{resolution}

%------------------------------------------------------------------------------%
\endinput

